\clearpage
\fancyhead[L]{\small GERLING.}
\hypertarget{native-547-lower}{}\label{native-547-lower}
\begin{center}\small Göttingische gelehrte Anzeigen. Stück 190. Seite 1889..1891. 1812 November 28.\\[3pt]\rule{.72\linewidth}{.3pt}\end{center}

\emph{Methodi proiectionis orthographicae usum ad calculos parallacticos facilitandos explicavit, simulque eclipsin solarem die VII. Sept. \textup{1820} apparituram hoc modo tractatam mappaque geographica illustratam tamquam exemplum proposuit} \textsc{Chr. Lud. Gerling}. 46 Seiten in Quart. Göttingen. Gedruckt bei \textsc{I. C. Baier}.

Zur Bestimmung der Umstände der Erscheinungen einer Sternbedeckung oder Sonnenfinsterniss für alle Punkte der Erde, wo dieselbe sichtbar ist, bedienen sich bekanntlich die Astronomen der orthographischen Projection auf eine Ebene, auf welcher die gerade Linie vom Mittelpunkte des bedeckten oder verfinsterten Gestirns senkrecht ist. Dies Verfahren gibt eine leichte und klare Vorstellung von den verschiedenen Punkten des Raumes, mit denen man es hierbei zu thun hat, und bahnt so einen bequemen Weg zur Auflösung der dabei vorkommenden Aufgaben, man möge sich der Zeichnung, oder des schärfsten aller Instrumente des Calculs bedienen wollen. Die vorliegende kleine Abhandlung, welche \SourceLeaf{548}von ihrem Verfasser --- ehemals unser gelehrter Mitbürger, jetzt Lehrer der Mathematik am königl. Lyceum in Cassel --- der hiesigen philosophischen Facultät bei Gelegenheit seiner Promotion als Probeschrift eingereicht wurde, gibt zu dieser Methode eine, gründliche Kenntnisse zeigende, fassliche Anleitung. Die Entwickelungen sind fast durchgehends analytisch ausgeführt, so dass der Verfasser gar nicht einmal eine Figur beizufügen nöthig gefunden hat. Dass dies ohne Schaden der Deutlichkeit unterlassen werden konnte, so dass jeder nur einigermassen Geübte den Schlüssen des Verfassers leicht folgen kann, ist ein Beweis der innern eigenthümlichen Klarheit der analytischen Methode bei geometrischen Aufgaben, wenn sie auf eine zweckmässige Art angewandt wird.

Die Sonnenfinsterniss vom 7. September 1820, auf welche, als Beispiel, der Verfasser die Methode anwendet, wird für unsere Gegend von Europa auf lange Zeit eine der merkwürdigsten sein. Die Oerter, für welche sie central ist, liegen in einer Linie, welche im nördlichsten Theile von Amerika anhebt, nördlich über Grönland und Island weg durch die Nordsee in der Richtung von Emden bis Triest durch Deutschland geht, sich dann durch das Adriatische Meer seiner Länge nach, über Morea und die nordlichste Spitze des rothen Meeres, zieht, und beim Untergange der Sonne in Arabien sich endigt. Für diese ganze Linie, und eine ziemlich breite Zone oberhalb und unterhalb derselben, ist die Finsterniss ringförmig, auch bei uns in Göttingen, wo der Ring, nach des Verfassers Rechnung, von 2 Uhr 40 Min. 41 Sec. bis 2 Uhr 46 Min. 0 Sec. wahrer Zeit dauern wird. Auch für Bremen, Seeberg, Berlin, Wien und Mannheim hat der Verfasser die Hauptumstände der Finsterniss hier mitgetheilt, und für eine noch grössere Anzahl von Puncten verspricht er, sie an einem andern Orte nachzuholen.
